% Generated draft: requires mathematical and source review.
\documentclass[11pt]{article}
\usepackage[margin=25mm]{geometry}
\usepackage{fontspec}
\setmainfont{Latin Modern Roman}
\usepackage{amsmath,amssymb,mathtools,unicode-math}
\setmathfont{Latin Modern Math}
\usepackage{enumitem,needspace,xurl,hyperref}
\setlength{\parindent}{0pt}
\setlength{\parskip}{5pt}
\setlength{\emergencystretch}{3em}
\newcommand{\N}{\mathbb{N}}
\newcommand{\Z}{\mathbb{Z}}
\newcommand{\Q}{\mathbb{Q}}
\newcommand{\R}{\mathbb{R}}
\newcommand{\C}{\mathbb{C}}
\newcommand{\F}{\mathbb{F}}
\newcommand{\A}{\mathbb{A}}
\newcommand{\PP}{\mathbb{P}}
\newcommand{\E}{\mathbb{E}}
\newcommand{\eps}{\varepsilon}
\newcommand{\dd}{\,\mathrm{d}}
\begin{document}
\section*{UnsolvedMath ID: 700840}
\textbf{Source problem code:} catalogue entry
\subsection{Definitions and mathematical statement}
Let \(R\) be a commutative ring with identity, and let \(G\) be a finite group. The group ring \(RG\) consists of formal sums \(\sum_{g\in G}r_g g\), with coefficients in \(R\), with coefficientwise addition and multiplication induced by the multiplication in \(G\). An \(R\)-representation of \(G\) on a unital \(R\)-module \(M\) is a homomorphism \(G\to\operatorname{Aut}_R(M)\), equivalently a unital left \(RG\)-module structure on \(M\).

The source writes \(a(RG)\) for the ring of \(R\)-representations of \(G\). It does not specify the admissible underlying \(R\)-modules or the relations used to construct this ring. Thus its construction remains an unresolved qualification: in particular, the supplied record does not say whether representations must be finite free, finitely generated projective, or otherwise restricted, or whether the additive relations arise from direct sums or short exact sequences.

An element \(e\) of a unital ring \(A\) is idempotent if \(e^2=e\); it is nontrivial if \(e\ne 0_A\) and \(e\ne 1_A\). Subject to clarification of the source's intended construction of \(a(RG)\), the requested assertion is: for every commutative ring \(R\) with identity and every finite group \(G\),
\[
\forall e\in a(RG),\qquad e^2=e\ \Longrightarrow\ e\in\{0_{a(RG)},1_{a(RG)}\}.
\]
No additional hypotheses on \(R\) are stated in the supplied record.
\subsection{Short English statement}
The problem asks whether the ring of representations of every finite group over every commutative ring with identity has only the idempotents \(0\) and \(1\). The precise construction intended by \(a(RG)\) requires clarification.
\subsection{Sources}
\begin{itemize}
\item Local UnsolvedMath registry, permanent ID \(700840\), exact source code \texttt{catalogue entry}, title \emph{Kourovka Notebook Problem 12.21}. The supplied record labels the problem \texttt{open}; this label is supplied registry information, not an independently verified mathematical or literature assessment. The record was updated on 2026-10-16.
\item \emph{The Kourovka Notebook (Unsolved Problems in Group Theory)}, problem 12.21, attributed in the supplied record to the 12th issue (1992) and the 21st-issue edition (2026), \texttt{arXiv:1401.0300v48}. Proposed by P. M. Gudivok and V. P. Rud'ko in 1992. These bibliographic details and the problem wording are supplied attribution; the cited edition has not been independently consulted.
\end{itemize}
\end{document}
